\chapter{LITERATURE REVIEW}
The scientific foundation of bone conduction as an auditory pathway was established by von Bek\'{e}sy \cite{bekesy1960}, whose experimental investigations demonstrated that mechanical vibrations applied to the skull can directly stimulate the fluid filled cochlea, bypassing the outer and middle ear entirely. This work underpinned the development of bone anchored hearing aids and, eventually, bone conduction communication headsets for occupational and military use.\\[10pt]
The frequency response characteristics of bone conducted voice were systematically studied by P\"{o}rschmann \cite{porschmann2000}, who performed side-by-side comparative recordings of the same speakers using both bone conduction and air conduction microphones. The study quantified a progressive high-frequency roll-off in the bone conduction channel above approximately 1--2~kHz, concluding that the skull acts as a low-pass mechanical filter for self-generated voice vibrations. This spectral distortion, if uncorrected, reduces consonant recognition rates and thus reduces speech intelligibility even in the absence of any external noise.\\[10pt]
Methods for combining air and bone conduction microphone signals for speech enhancement were pioneered by Liu~et~al. \cite{liu2004}, who introduced a speaker independent ``direct filtering'' technique that constructs an inverse filter from the spectral ratio of simultaneously recorded air conduction and bone conduction speech. Their experiments demonstrated that this approach effectively removes non-stationary environmental noise and significantly improves automatic speech recognition accuracy, directly informing the spectral inverse filter approach used for FIR equalization in the present project.\\[10pt]
The SSM2167 microphone preamplifier IC, manufactured by Analog Devices, is a single-chip solution that integrates a low noise preamplifier, a voltage controlled amplifier (VCA) for dynamic compression with programmable compression ratio (up to 10:1), and a downward expander noise gate that physically attenuates the output when the input signal falls below a programmable threshold. As specified in the SSM2167 datasheet \cite{ssm2167}, the device achieves less than 0.2\% THD+N across a 20~kHz bandwidth and supports compression ratios from 2:1 to 10:1, with the noise gate preventing amplification of background hum or low-level ambient noise when no speech is present, all without any digital processing. This property makes it the cornerstone of the analog noise suppression strategy in the proposed system.\\[10pt]
Active filter design using operational amplifiers for audio-band bandpass filtering is a well-established technique. Sedra and Smith \cite{sedra2014} provide the theoretical basis for Sallen Key and multiple feedback topologies using general purpose op-amps, demonstrating that second and fourth order Butterworth bandpass filters constructed from the NE5532 offer low noise, sufficient bandwidth, and excellent stopband attenuation for voice-band audio applications. Components for these filter networks are universally available and inexpensive.\\[10pt]
The practical adoption of TinyML on STM32 microcontrollers was substantially enabled by the STM32Cube.AI tool from STMicroelectronics, which compiles Keras and TensorFlow Lite models into optimized C code deployable on Cortex-M4 targets. David~et~al. \cite{david2021} benchmarked a range of quantized neural network models on the STM32F4 series and showed that Dense Neural Networks (DNNs) and small Recurrent Neural Networks (RNNs) with fewer than 50~KB of parameters can run real-time inference on 20~ms audio frames at 16~kHz within the available compute budget of the STM32F401. This directly informs the feasibility and scope of the TinyML transient detection component in the proposed project.\\[10pt]
Real-time transient noise suppression distinct from steady-state noise suppression has received comparatively less attention in the literature. Zhou~et~al. \cite{zhou2020} demonstrated that a bone-conduction-assisted dual-microphone speech enhancement system can achieve robust speech detection and noise suppression in real time, using the bone-conduction channel as a reliable voice activity reference to control the adaptive noise cancellation loop, outperforming single-microphone approaches when impulse noise and non-stationary background noise co-occur. Their architecture which uses the noise-immune bone-conduction signal to gate and guide a digital enhancement stage is directly analogous to the VAD and TinyML-gated limiter design in the present project.\\[10pt]
Frequency equalization of bone conduction speech through reconstruction filtering was formalised by Shimamura and Tamiya \cite{shimamura2006}, who designed a linear-phase finite impulse response filter using the long-term spectral characteristics of simultaneously recorded bone-conduction and normal air-conduction speech from the same speaker. Their work demonstrated that such a reconstruction filter can meaningfully restore the high-frequency content attenuated by the bone-conduction transduction path, and that the filter's properties can be investigated and validated through spectral analysis of the reconstructed speech signal. This directly supports the offline FIR filter design methodology adopted in the proposed project.\\[10pt]
The broader physiological and engineering basis for bone-conduction hearing and communication systems was comprehensively reviewed by Stenfelt and Goode \cite{stenfelt2005}, who identified five distinct mechanisms by which skull vibration excites the cochlea and characterised the frequency-dependent contributions of each. Their review established that bone-conduction hearing is highly effective across the speech-frequency range and provided the quantitative physiological foundation for understanding why bone-conduction microphones achieve high voice-to-noise selectivity in the sub-3~kHz range while exhibiting the well-known high-frequency roll-off that the FIR equalization stage in the present project is designed to correct.